\chapter{Algorithmique, listes et calcul fiable}
\section*{Exercice 24.1 -- Listes et indices}
Que produit \texttt{[k*k for k in range(1, 5)]} ? Que valent les éléments d'indices $0$ et $-1$ ?

\section*{Exercice 24.2 -- Copie ou alias}
Après \texttt{a = [1, 2]}, \texttt{b = a}, \texttt{b.append(3)}, que vaut \texttt{a} ? Comment éviter cet effet ?

\section*{Exercice 24.3 -- Invariant}
Justifier l'algorithme qui initialise $s=0$, puis ajoute successivement $1,2,\ldots,n$.

\section*{Exercice 24.4 -- Syracuse et borne}
Écrire le début de Syracuse à partir de $6$. Expliquer pourquoi un programme général doit prévoir une borne d'itérations.

\section*{Exercice 24.5 -- Différences finies}
En arithmétique exacte, quel résultat donne $[(a+h)^2-a^2]/h$ ? Pourquoi un très petit $h$ peut-il poser problème en flottants ?

\section*{Exercice 24.6 -- Fréquences de lettres}
En ignorant espaces et accents, calculer les fréquences des lettres de « été ». Quelle précaution faut-il pour comparer deux langues ?
